% GENERATED FILE. Edit canonical lessons, then run npm run generate:books.
% canonical-source-hash: fnv1a64:7d3a3f5f422f2bd7
% canonical-lessons: BN-C526-setar, BN-C526-dhol, BN-C526-harmoniyam, BN-C526-behala, BN-C526-tungtang-kora

\chapter{Instruments}
\label{ch:bn-instruments}
\hlchaptermodality{526}

\begin{chapteropening}
\textbf{By the end of this chapter:} \emph{I can name some instruments and say how they sound.}

Complete the last lesson of chapter 526: I can name some instruments and say how they sound.
\end{chapteropening}

\section[\texorpdfstring{\bn{সেতার} (setār)}{সেতার (setār)}]{\bn{সেতার} (setār) — a sitar}

\label{lesson:BN-C526-setar}



\begin{quote}
Before the new one: say the Bengali for to resound, then the Bengali for to whistle.
\end{quote}

\subsection*{You'll want to know: \bn{সেতার}}
\textbf{\bn{সেতার}} — \emph{setār} — \textquotedblleft{}a sitar\textquotedblright{}.

\begin{quote}
\bn{আমার} \bn{কাকা} \bn{সেতার} \bn{বাজান।}
\end{quote}

\begin{quote}
\emph{Āmār kākā setār bājān.}
\end{quote}

\begin{quote}
My uncle plays the sitar.
\end{quote}

\begin{cousinweb}
\textbf{\bn{সেতার}} (\emph{setār}) is usually traced to Persian \emph{sitār}, \textquotedblleft{}three strings\textquotedblright{}: \emph{si}, \textquotedblleft{}three\textquotedblright{}, and \emph{tār}, \textquotedblleft{}string\textquotedblright{}.
\end{cousinweb}

\begin{grammarlens}[title={What you've built}]
One more word for describing what you lived and felt.
\end{grammarlens}

\subsection*{Your turn}
\begin{itemize}
\raggedright
  \item \emph{Say it:} \emph{setār}
  \item \emph{Say it:} \emph{setār}, once more
  \item \emph{Say it:} \emph{Āmār kākā setār bājān.} — My uncle plays the sitar.
  \item \emph{Recall:} say the Bengali for to resound, then the Bengali for to whistle, then say \emph{setār} again
\end{itemize}

\begin{tcolorbox}[breakable,colback=teal!4,colframe=teal!35!black,title={Before you move on}]
What does \bn{সেতার} mean? (\textquotedblleft{}A sitar\textquotedblright{}.) Say it once more.
\end{tcolorbox}
\section[\texorpdfstring{\bn{ঢোল} (ḍhol)}{ঢোল (ḍhol)}]{\bn{ঢোল} (ḍhol) — a drum; a dhol (a two-headed drum)}

\label{lesson:BN-C526-dhol}



\begin{quote}
Before the new one: say the Bengali for to whistle, then the Bengali for a sitar.
\end{quote}

\subsection*{You'll want to know: \bn{ঢোল}}
\textbf{\bn{ঢোল}} — \emph{ḍhol} — \textquotedblleft{}a drum; a dhol (a two-headed drum)\textquotedblright{}.

\begin{quote}
\bn{মেলায়} \bn{একজন} \bn{ঢোল} \bn{বাজাচ্ছিল।}
\end{quote}

\begin{quote}
\emph{Melāy ekjôn ḍhol bājāchchhilô.}
\end{quote}

\begin{quote}
Someone was playing a drum at the fair.
\end{quote}

\begin{cousinweb}
\textbf{\bn{ঢোল}} (\emph{ḍhol}) is a tadbhava: CDIAL lists Bengali \emph{ḍhol}, \textquotedblleft{}drum\textquotedblright{}, under \emph{ḍhōla} (CDIAL 5608).
\end{cousinweb}

\begin{grammarlens}[title={What you've built}]
One more word for describing what you lived and felt.
\end{grammarlens}

\subsection*{Your turn}
\begin{itemize}
\raggedright
  \item \emph{Say it:} \emph{ḍhol}
  \item \emph{Say it:} \emph{ḍhol}, once more
  \item \emph{Say it:} \emph{Melāy ekjôn ḍhol bājāchchhilô.} — Someone was playing a drum at the fair.
  \item \emph{Recall:} say the Bengali for to whistle, then the Bengali for a sitar, then say \emph{ḍhol} again
\end{itemize}

\begin{tcolorbox}[breakable,colback=teal!4,colframe=teal!35!black,title={Before you move on}]
What does \bn{ঢোল} mean? (\textquotedblleft{}A drum; a dhol (a two-headed drum)\textquotedblright{}.) Say it once more.
\end{tcolorbox}
\section[\texorpdfstring{\bn{হারমোনিয়াম} (hārmoniyām)}{হারমোনিয়াম (hārmoniyām)}]{\bn{হারমোনিয়াম} (hārmoniyām) — a harmonium}

\label{lesson:BN-C526-harmoniyam}



\begin{quote}
Before the new one: say the Bengali for a sitar, then the Bengali for a drum.
\end{quote}

\subsection*{You'll want to know: \bn{হারমোনিয়াম}}
\textbf{\bn{হারমোনিয়াম}} — \emph{hārmoniyām} — \textquotedblleft{}a harmonium\textquotedblright{}.

\begin{quote}
\bn{দিদিমার} \bn{একটা} \bn{পুরনো} \bn{হারমোনিয়াম} \bn{আছে।}
\end{quote}

\begin{quote}
\emph{Didimār ekṭā purôno hārmoniyām āchhe.}
\end{quote}

\begin{quote}
My grandmother has an old harmonium.
\end{quote}

\begin{cousinweb}
\textbf{\bn{হারমোনিয়াম}} (\emph{hārmoniyām}) is borrowed from English \emph{harmonium}.
\end{cousinweb}

\begin{grammarlens}[title={What you've built}]
One more word for describing what you lived and felt.
\end{grammarlens}

\subsection*{Your turn}
\begin{itemize}
\raggedright
  \item \emph{Say it:} \emph{hārmoniyām}
  \item \emph{Say it:} \emph{hārmoniyām}, once more
  \item \emph{Say it:} \emph{Didimār ekṭā purôno hārmoniyām āchhe.} — My grandmother has an old harmonium.
  \item \emph{Recall:} say the Bengali for a sitar, then the Bengali for a drum, then say \emph{hārmoniyām} again
\end{itemize}

\begin{tcolorbox}[breakable,colback=teal!4,colframe=teal!35!black,title={Before you move on}]
What does \bn{হারমোনিয়াম} mean? (\textquotedblleft{}A harmonium\textquotedblright{}.) Say it once more.
\end{tcolorbox}
\section[\texorpdfstring{\bn{বেহালা} (behālā)}{বেহালা (behālā)}]{\bn{বেহালা} (behālā) — a violin}

\label{lesson:BN-C526-behala}



\begin{quote}
Before the new one: say the Bengali for a drum, then the Bengali for a harmonium.
\end{quote}

\subsection*{You'll want to know: \bn{বেহালা}}
\textbf{\bn{বেহালা}} — \emph{behālā} — \textquotedblleft{}a violin\textquotedblright{}.

\begin{quote}
\bn{মেয়েটা} \bn{স্কুলে} \bn{বেহালা} \bn{বাজাতে} \bn{শেখে।}
\end{quote}

\begin{quote}
\emph{Meyeṭā skule behālā bājāte shekhe.}
\end{quote}

\begin{quote}
The girl is learning to play the violin at school.
\end{quote}

\begin{cousinweb}
\textbf{\bn{বেহালা}} (\emph{behālā}) is usually traced to Portuguese \emph{viola}, a stringed instrument.
\end{cousinweb}

\begin{grammarlens}[title={What you've built}]
One more word for describing what you lived and felt.
\end{grammarlens}

\subsection*{Your turn}
\begin{itemize}
\raggedright
  \item \emph{Say it:} \emph{behālā}
  \item \emph{Say it:} \emph{behālā}, once more
  \item \emph{Say it:} \emph{Meyeṭā skule behālā bājāte shekhe.} — The girl is learning to play the violin at school.
  \item \emph{Recall:} say the Bengali for a drum, then the Bengali for a harmonium, then say \emph{behālā} again
\end{itemize}

\begin{tcolorbox}[breakable,colback=teal!4,colframe=teal!35!black,title={Before you move on}]
What does \bn{বেহালা} mean? (\textquotedblleft{}A violin\textquotedblright{}.) Say it once more.
\end{tcolorbox}
\section[\texorpdfstring{\bn{টুংটাং} \bn{করা} (ṭungṭāng kôrā)}{টুংটাং করা (ṭungṭāng kôrā)}]{\bn{টুংটাং} \bn{করা} (ṭungṭāng kôrā) — to strum; to tinkle}

\label{lesson:BN-C526-tungtang-kora}



\begin{quote}
Before the new one: say the Bengali for a harmonium, then the Bengali for a violin.
\end{quote}

\subsection*{You'll want to know: \bn{টুংটাং} \bn{করা}}
\textbf{\bn{টুংটাং} \bn{করা}} — \emph{ṭungṭāng kôrā} — \textquotedblleft{}to strum; to tinkle\textquotedblright{}.

\begin{quote}
\bn{সন্ধ্যায়} \bn{মেয়েটা} \bn{গিটারে} \bn{টুংটাং} \bn{করছিল।}
\end{quote}

\begin{quote}
\emph{Sôndhyāy meyeṭā giṭāre ṭungṭāng kôrchhilô.}
\end{quote}

\begin{quote}
In the evening the girl was strumming on the guitar.
\end{quote}

Past (he/she …): \bn{তিনি} \bn{সেতারে} \bn{টুংটাং} \bn{করলেন}, \emph{tini setāre ṭungṭāng kôrlen}.

\begin{cousinweb}
\textbf{\bn{টুংটাং} \bn{করা}} (\emph{ṭungṭāng kôrā}) joins \bn{টুংটাং}, an imitative word for the light sound of plucked strings or a small bell, to \bn{করা}, \textquotedblleft{}to do\textquotedblright{}. No dictionary we use records an older form, so its history is not settled.
\end{cousinweb}

\begin{grammarlens}[title={What you've built}]
One more word for describing what you lived and felt.
\end{grammarlens}

\subsection*{Your turn}
\begin{itemize}
\raggedright
  \item \emph{Say it:} \emph{ṭungṭāng kôrā}
  \item \emph{Say it:} \emph{ṭungṭāng kôrā}, once more
  \item \emph{Say it:} \emph{Sôndhyāy meyeṭā giṭāre ṭungṭāng kôrchhilô.} — In the evening the girl was strumming on the guitar.
  \item \emph{Recall:} say the Bengali for a harmonium, then the Bengali for a violin, then say \emph{ṭungṭāng kôrā} again
\end{itemize}

\begin{tcolorbox}[breakable,colback=teal!4,colframe=teal!35!black,title={Before you move on}]
What does \bn{টুংটাং} \bn{করা} mean? (\textquotedblleft{}To strum; to tinkle\textquotedblright{}.) Say it once more.
\end{tcolorbox}
